\chapter{Un neurone, plusieurs appareils}
% version complète: chapters/18_eetypes.tex

Au premier chapitre, nous avons trié les neurones selon la substance qu'ils libèrent. Un ingénieur les trierait autrement: selon l'appareil que chacun réalise. Et il se trouve que, sous une apparence identique, se cache toute une boîte de composants électroniques.

\section*{La boîte de composants}

Le seau percé du deuxième chapitre est un intégrateur à fuite: il additionne les signaux sur une courte fenêtre. Avec une fenêtre très petite, il devient un détecteur de coïncidences. Le neurone qui ne répond qu'aux changements et se tait aussitôt est un filtre qui laisse passer les variations. Le neurone qui travaille tout seul est un métronome, un oscillateur; si une entrée ajuste son tempo, c'est un oscillateur commandé en tension. Le neurone qui émet des bouffées de clics est un générateur de bouffées, qui écrit des \og traits\fg{}. Le neurone qui clique au rythme de l'onde sonore est un détecteur de phase.

\pencilfig{18}{\pencilart{18}}{Sous une apparence identique se cache toute une boîte de composants: résistance, condensateur, bobine, oscillateur.}

Il y a aussi des appareils dont nous n'avons pas encore parlé. Le \emph{résonateur} ressemble à une balançoire: il répond le mieux aux poussées d'un certain rythme, comme une balançoire qui ne prend de l'élan que si on la pousse au bon moment. Dans le neurone, le rôle du ressort est tenu par un courant lent qui s'oppose aux changements. L'\emph{intégrateur sans fuite} garde une valeur longtemps, comme un condensateur sans fuite: c'est ainsi que fonctionne le réseau qui maintient la position des yeux quand vous regardez de côté. Mais il faut pour cela une rétroaction ajustée presque parfaitement, et ce réglage coûte cher. La \emph{bascule} est un neurone qu'une brève poussée fait passer à un travail prolongé jusqu'à ce qu'on l'arrête, comme un interrupteur à verrouillage; chez les neurones qui commandent les muscles, ce mode est activé par la sérotonine.

\section*{Un appareil caméléon}

La grande découverte de ce tri: ce ne sont pas des composants différents, mais différents \emph{modes} d'un seul et même composant. Un même neurone peut être intégrateur le jour et générateur de bouffées la nuit, selon ce que lui disent les stations de radio. Un ingénieur verrait dans le cerveau un circuit logique reprogrammable: le circuit est unique, et ce sont les réglages qui changent la fonction des composants.

Que les neurones savent fonctionner dans tous ces modes, c'est mesuré. Comment le cerveau utilise le changement de mode pour calculer, c'est en grande partie une hypothèse.

\fullversion{18}
